\section{Discussions}

\label{sec:discussions}
As started previously, we used Nomic Embed data to enable direct comparison with existing models. While this choice grounds our results in a well-established setup, it also means that some observations may not generalize to different or stronger training configurations, an avenue we leave for future work.
For example, when performing the supervised fine-tuning phase on stronger data (NV-Retriever~\cite{moreira2024nv} on various datasets), adding prompts on top of a ColBERT model pre-trained without prompts did not degrade the results as it was the case for ColBERT-Zero. With stronger/longer fine-tuning, the alignment with pre-training is less required as the model has enough room to adapt. Likewise, the small remaining advantage of pre-training a ColBERT model, rather than relying solely on fine-tuning and KD, is likely influenced by the quality and scale of those subsequent phases. Given that knowledge distillation typically \emph{improves} results after supervised fine-tuning, the fact that a prior contrastive phase further enhances performance suggests that these gains are less about the specific objective and more about the \emph{scale of training}. This extra capacity allows the model more time to optimize within a multi-vector setting. Consequently, performing KD alone at a larger scale might yield similar, if not superior, results due to the higher quality of the distillation signal. The goal of this study, however, was to explore the conventional setup where training scale is inversely proportional to signal quality, reflecting the higher cost of generating high-quality labels. Thus, the results of each phase primarily highlight the impact of \emph{scale} within those phases rather than the specific training objective itself.
% We thus retrain the model starting from the dense model but append prompts during training and evaluation to ensure the best setup. We observe a strong gain (+1.25), confirming that alignment with the base model is very important. 

% we observe surprising results: the final model is actually better than the fully pre-trained ColBERT and even \texttt{GTE-ModernColBERT} that leverage the very strong gte-modernbert model. Given that we do not observe a big difference in the score by using longer lengths (both during training and evaluation), this factor cannot explain the difference alone.



% The first thing to note is that, the fully pre-trained ColBERT model almost matches the results of \texttt{GTE-ModernColBERT}. This is very strong results, because the base model used for the latter (GTE-ModernBERT) is a very performant model optimized on strong non-public data. The superiority of the GTE data is shown by the huge gap in BEIR results for the dense models (52.9 for Nomic data vs 55.33 for GTE data). 

% The first thing to note is that, the fully pre-trained ColBERT model almost matches the results of \texttt{GTE-ModernColBERT}. This is very strong results, because the base model used for the latter (GTE-ModernBERT) is a very performant model optimized on strong non-public data. The superiority of the GTE data is shown by the huge gap in BEIR results for the dense models (52.9 for Nomic data vs 55.33 for GTE data). Consequently, achieving comparable results highlights that running a full ColBERT pre-training before performing knowledge distillation is noticeably better than only performing KD on top of a pre-trained dense model. This is further confirmed by the results obtained by performing the KD step on top of the dense supervised model from Nomic, which only obtains a score of 52.58, lagging much behind the fully ColBERT-pre-trained model by 2 pts of nDCG@10. 
% Interestingly, it should be noted that the resulting ColBERT actually achieve a \textbf{worse} BEIR score than the dense model it is derived from. Although surprising, it should be noted that this was also the case of \texttt{GTE-ModernColBERT}, achieving worse results than its base model, gte-modernbert.

% However, in order to provide the fairest comparison possible, it is important to note that the dense models from Nomic are actually trained with \textit{prompts} (respectively "search\_query:" and "search\_document:"). Unlike approaches such as Promptriever~\cite{weller2025promptriever}, the goal of these prompts is less about specializing the models over different tasks and making it "promptable", but rather creating an asymmetric encoding. Although ColBERT already has an asymmetric encoding mechanism baked-in (through the [Q] and [D] identifiers appended to queries/documents, much like prompts), having a mismatch between the pre-trained model and the fine-tuning setup could hurt the performance. We thus retrain the model starting from the dense model but append prompts during training and evaluation to ensure the best setup. We observe a strong gain (+1.25), confirming that alignment with the base model is very important. 

% However, the model still lag behind the fully pre-trained model but by a much closer margin. An interesting question is whether we could close the gap by running a supervised phase in the ColBERT setting before the KD phase, making the training a bit heavier than only performing KD but still skipping the most expensive unsupervised phase. When running both the supervised and knowledge distillation phase starting from the unsupervised dense model (using prompts to be aligned with the base model), we observe surprising results: the final model is actually better than the fully pre-trained ColBERT and even \texttt{GTE-ModernColBERT} that leverage the very strong gte-modernbert model. Given that we do not observe a big difference in the score by using longer lengths (both during training and evaluation), this factor cannot explain the difference alone.



% When running the supervised and KD phases with prompts starting from the unsupervised ColBERT-Zero that has \textit{not} been trained with prompts, we observe no improvement (actually, a slight degradation) over not using prompts.
% This is most likely explained by the mismatch between the supervised phase and the two fine-tuning steps. To confirm this hypothesis, we re-run a whole ColBERT pre-training, integrating the prompts right from the start. As expected, we observe an improvement, with a final model outperforming not only the best ColBERT model obtained using a dense base model (and so \texttt{GTE-ModernColBERT}) but also gte-modernbert itself. The significance of this result is best appreciated in context: at the time of writing the model is the top-1 model on the MTEB BEIR leaderboard for model <150M parameters (top-4 for models <300M).
% To give some perspective, at the time of writing, the model has a lead of more than 2.5 nDCG@10 over ModernBERT-embed-supervised, which translate to a top-4 on leaderboard for model <300M (resp. top-1 for model <150M) against a rank 23 (resp. top-11), \textbf{despite using the same data}.
% \section{Discussions}
% We found that only a small knowledge distillation step for ColBERT training on top of a dense model is not enough to get the best performance possible, as illustrated by the improvement of more than 1 nDCG@10 when performing the supervised phase before KD. Running a more thorough supervised phase before largely improves the performance. The full pre-training of ColBERT still seems to have an edge but it is very contained (especially given the potential noise of BEIR). Moreover, this gap will surely be even smaller the better (higher quality/bigger quantity) the supervised phase. Likewise, it seems that running the pre-training in the ColBERT setting yields a small advantage over running only the fine-tuning. However, it is probable that, the heavier the fine-tuning, the less likely this is to matter.

% We also observe that prompts had a large impact on the results. However, it is hard to properly disentangle all the possible reasons explaining the differences in performance as there are different effect at play. Aligning the fine-tuning with the pre-trained model seems to be beneficial, as noted by the gain brought by the addition of the prompts during the KD phase, even without changing the lengths at all. 
%     \item The additional tokens can serve as query (/document) expansion. This mechanism has shown very useful in the early variant of ColBERT~\cite{khattab2020colbert}. Typically, these tokens serve as "placeholder" tokens that does not correspond to a specific semantic token but can be used by the model to store global information about the sequence. Despite being very helpful, this mechanism is not used in modern models leveraging Flash Attention\footnote{Query expansion worked by used PAD tokens that were not attended by other tokens in the sequence. In the original (not optimized) attention implementations, all the attentions scores were computed and the contributions of masked tokens were zeroed out at the end so that the representations of the tokens were only computed w.r.t unmasked tokens. However, the representations of the masked tokens w.r.t all the others tokens were still computed. Those are usually thrown away using attention mask afterwards, but in the case of ColBERT, those were used as query expansion. With more recent and optimized implementations of attention such as Flash Attention, the embedding of masked tokens are zeros/NaNs, preventing their usage as query expansion.}~\cite{dao2022flashattention} such as ModernBERT~\cite{DBLP:conf/acl/WarneTGBLA25}, the base model used by all models in our experiments.

% The key results to extract are the following:
% \begin{itemize}
%     \item Only a small knowledge distillation step for ColBERT training on top of a dense model model is not enough to get the best performance possible, as illustrated by the improvement of more than 1 nDCG@10 when performing the supervised phase before KD.
%     \item Running a more thorough supervised phase before largely improves the performance. The full pre-training of ColBERT still seems to have an edge but it is very contained (especially given the potential noise of BEIR). Moreover, this gap will surely be even smaller the better (higher quality/bigger quantity) the supervised phase.
  
%     % \item Likewise, it seems that running the pre-training in the ColBERT setting yields a small advantage over running only the fine-tuning. However, it is likely that, the heavier the fine-tuning, the less likely to matter much. The goal of this paper was originally to only use the nomic data for (un)supervised phases, to re-use the existing models and make comparison.
%     % \item Aligning the fine-tuning with the pre-trained model seems to be beneficial, as noted by the gain brought by the addition of the prompts during the KD phase, even without changing the lengths at all. 
% \end{itemize}


% About the effect of prompts: It is hard to disentangle all the possible reasons explaining the difference in performances. There are different effect at play here: 
% \begin{itemize}
    
%     \item Aligning the fine-tuning with the pre-trained model seems to be beneficial, as noted by the gain brought by the addition of the prompts during the KD phase, even without changing the lengths at all. 
%     \item However, this alignment does not explain all the gains, as we observe that full pre-training with and without prompts have widely different scores (54.61 vs 55.43).
%     \item Longer lengths could prevent queries/documents in the evaluations to be truncated during the evaluation. However, simply running the evaluation with unprompted models but longer lengths, we did not observe any change.
%     Likewise, even when with the additional prompts, increasing the lengths had marginal effects, most likely preventing the truncation due to the additional tokens brought by prompts.
%     \item asymmetric encoding through prompts has shown helpful and is widely used~\cite{wang2022text,nussbaum2025nomic}. Although ColBERT models already have specific query/document markers, more tokens could be more explicit.
%     \item The additional tokens can serve as query (/document) expansion. This mechanism has shown very useful in the early variant of ColBERT~\cite{khattab2020colbert}. Typically, these tokens serve as "placeholder" tokens that does not correspond to a specific semantic token but can be used by the model to store global information about the sequence. Despite being very helpful, this mechanism is not used in modern models leveraging Flash Attention\footnote{Query expansion worked by used PAD tokens that were not attended by other tokens in the sequence. In the original (not optimized) attention implementations, all the attentions scores were computed and the contributions of masked tokens were zeroed out at the end so that the representations of the tokens were only computed w.r.t unmasked tokens. However, the representations of the masked tokens w.r.t all the others tokens were still computed. Those are usually thrown away using attention mask afterwards, but in the case of ColBERT, those were used as query expansion. With more recent and optimized implementations of attention such as Flash Attention, the embedding of masked tokens are zeros/NaNs, preventing their usage as query expansion.}~\cite{dao2022flashattention} such as ModernBERT~\cite{DBLP:conf/acl/WarneTGBLA25}, the base model used by all models in our experiments.
% \end{itemize}








% \subsection{Knowledge distillation on pre-trained multi-vector and single-vector}
% The first interesting comparison is models obtained after knowledge distillation on top of the dense and multi-vector pre-trained models. The results (54.61 for the multi-vector pre-trained model vs 53.83 for the adapted dense) clearly confirm that pre-training in a multi-vector setup yields better multi-vector model after knowledge distillation than performing in the single-vector setting. Interestingly, the fully pre-trained ColBERT model almost matches the results of \texttt{GTE-ModernColBERT}. This is very impressive, because the base model used for the latter (GTE-ModernBERT) is a very performant model optimized on strong non-public data - the superiority of the GTE data is shown by the huge gap in BEIR results for the dense models (52.9 for Nomic data vs 55.33 for GTE data). Consequently, achieving comparable results highlights that running a full ColBERT pre-training before performing knowledge distillation is noticeably better than only performing KD on top of a pre-trained dense model.

% Even more interestingly, we can see that the actual BEIR results \textit{degrades} when performing the KD process w.r.t the dense base model results. We originally observed this when training \texttt{GTE-ModernColBERT} but we considered it could be due to the base model being strongly fitted to BEIR. However, this happens again while ModernBERT-base-embed is not as fitted to BEIR and thus seems to be a more general behavior. When performing the knowledge distillation on the ColBERT-pre-trained model, on the other hand, the BEIR performance strongly improves as expected.
% One possible explanation for this difference is observing another phenomenon: when performing KD on single-vector models, the BEIR score degrades while the NanoBEIR score steadily increases.
% Although NanoBEIR sometimes can be slightly uncorrelated with BEIR, this can also be observed when performing KD on both \texttt{GTE-ModernColBERT} and ModernBERT-embed.
% One possible explanation is that, while the BEIR evaluation uses the PLAID index because the datasets are too large for an exhaustive search, the NanoBEIR evaluation uses an exact search. Although the small KD step enables the model to work in a multi-vector setting, the resulting embedding space may not be sufficiently trained to work perfectly with PLAID. On the other hand, the large pre-training creates a space that is carefully optimized for late-interaction. However, these results only do not allow us to draw a clear conclusion on whether this hypothesis is correct, and it could be due to other confounding factors (e.g, architecture as both models are ModernBERT-based, the KD data as both use the same data…)

% \section{Contrastive Steps Comparison}
% Previous results highlights the benefits of performing the pre-training step in a multi-vector setting. One interesting question is whether this comes solely from better alignment during the KD phase or if the pre-trained models competitive are stronger than dense models before this phase. We thus compare single and multi-vector models after each phase.

% First, we can see that after the unsupervised phase, the gap is very large (51.70 for multi-vector vs 47.05 for single-vector), indicating that the multi-vector setting is able to better leverage the large unsupervised datasets and/or converge way faster than a single-vector model. Surprisingly, after the supervised phase, the dense model closes this large gap and ends up outperforming the late interaction model.
% This phenomenon could be explained by the multi-vector model converging earlier during the unsupervised phase and thus benefiting less from the supervised phase. Nevertheless, it is evident that ColBERT pre-training exhibits a high degree of competitiveness with dense pre-training. We also release the unsupervised and supervised checkpoints to allow further explorations
